% Generated draft: requires mathematical and source review.
\documentclass[11pt]{article}
\usepackage[margin=25mm]{geometry}
\usepackage{fontspec}
\setmainfont{Latin Modern Roman}
\usepackage{amsmath,amssymb,mathtools,unicode-math}
\setmathfont{Latin Modern Math}
\usepackage{enumitem,needspace,xurl,hyperref}
\setlength{\parindent}{0pt}
\setlength{\parskip}{5pt}
\setlength{\emergencystretch}{3em}
\newcommand{\N}{\mathbb{N}}
\newcommand{\Z}{\mathbb{Z}}
\newcommand{\Q}{\mathbb{Q}}
\newcommand{\R}{\mathbb{R}}
\newcommand{\C}{\mathbb{C}}
\newcommand{\F}{\mathbb{F}}
\newcommand{\A}{\mathbb{A}}
\newcommand{\PP}{\mathbb{P}}
\newcommand{\E}{\mathbb{E}}
\newcommand{\eps}{\varepsilon}
\newcommand{\dd}{\,\mathrm{d}}
\begin{document}
\section*{UnsolvedMath ID: 30004145}
\textbf{Source problem code:} catalogue entry
\subsection{Definitions and mathematical statement}
Under the standard convention, a plane projective \(M\)-curve of degree \(3\) is a nonsingular real plane cubic \(X\subset\PP^2_{\R}\) whose real locus has the maximum number of connected components. Its genus is \(g=(3-1)(3-2)/2=1\), so write
\[
X(\R)=C_0\sqcup C_1.
\]
A real divisor on \(X\) is a finite formal sum \(A=\sum_{p\in X(\C)}m_p[p]\), with integer coefficients invariant under complex conjugation. Its degree is \(\sum_p m_p\). The supplied problem assumes that \(A\) is a transverse real divisor of degree \(9\), but does not define transversality or specify whether every support point must be real. In particular, invariance under conjugation alone does not imply that all support points are real.

\textbf{Requested determination.} For every such \(X\) and admissible \(A\), determine whether
\[
\pi_1\bigl(\operatorname{Emb}(A,X),e_A\bigr)=\{1\},
\]
and determine how the answer depends on the distribution of \(A\) among \(C_0\) and \(C_1\). Here \(e_A\) denotes the reference embedding representing the given divisor, provided that this embedding belongs to the intended space. If \(A\) consists of nine distinct real points, its component counts are
\[
n_i=\#\bigl(\operatorname{supp}(A)\cap C_i\bigr),\qquad n_0+n_1=9.
\]
This count description is conditional on total reality and reducedness; neither is explicitly defined in the supplied record.

\textbf{Unresolved formulation.} The record does not define \(\operatorname{Emb}(A,X)\): its objects, topology, labeling conventions, admissibility conditions, and any restrictions arising from divisors or plane intersections are unspecified. It also does not specify the relevant path component or basepoint. Thus the displayed fundamental-group question records the intended target, but a uniquely defined mathematical problem requires these missing definitions. No assertion of triviality or classification is made here.
\subsection{Short English statement}
For a real plane \(M\)-cubic and a transverse real divisor of degree \(9\), determine whether the intended embedding space has trivial fundamental group and how this depends on the divisor's placement on the two real components. The embedding space and transversality are undefined in the supplied record, so the precise question remains conditional on their clarification.
\subsection{Sources}
\begin{itemize}
\item Supplied local UnsolvedMath registry record, permanent ID \(30004145\), title \emph{Embedding Spaces of Divisors on Real Cubics}; collection: Oberwolfach Reports Open Problems; exact source problem code \texttt{catalogue entry}. No individual proposer is supplied; the proposed year is recorded as \(2019\).
\item Supplied report citation: \emph{Tropical Geometry: new directions} (2020), Oberwolfach Reports, pp.~1245--1308, DOI \texttt{10.4171/owr/2019/21}. \url{https://doi.org/10.4171/owr/2019/21}. The background describes Oberwolfach Reports as documenting mathematical workshops and their open-problem discussions.
\item Supplied literature source: \emph{Tropical geometry problem report}. \url{https://ems.press/content/serial-article-files/46800}.
\item Literature assessment dated exactly \texttt{2026-08-22}: the supplied curation labels the problem \texttt{partially\_solved} and states that a trace-homomorphism theorem and explicit low-degree cases resolve substantial instances, while a general triviality proof or complete description when the group is nontrivial remains unavailable. This is supplied dated curation, not an independently verified literature assessment. The record itself requests independent verification of time-sensitive status and cited claims.
\item The registry labels the statement \texttt{corrected\_verified} and reports manual semantic, typographic, and LaTeX review asserting self-containedness and preservation of the target. These are supplied editorial claims; they do not resolve the missing definitions identified above. No external source was consulted for this draft.
\end{itemize}
\end{document}
